\section{A Selmer presentation at a point with two blocks}

Let $S$ be any native $\mathrm{Ctx}(\ell)$, $A=k[[t]]$,
$F=k((t))$, $R=S.Bh$ with maximal ideal $\mathfrak m$,
$O'=S.OO$, and $r:R\twoheadrightarrow F$ its residue map.
Write $\rho$ for the point representation pulled back along
$S.h$, $g=S.k.g$, and $v$ for the representation in $S.U$.
Its finite commutative coefficient algebra is $R\oplus N$,
with $N^2\subset\mathfrak m\subset R$. The span of $v$ over
$R$ has diagonal blocks over $R$ and both opposite blocks over
$N$. Let $I$ be the ideal of $R$ generated by products of two
elements of $N$, and set
\[
 J=\mathfrak m_{O'}R+\mathfrak m^2+I .
\]
Write $K_1(D)$, $K_2$ and $K_3$ for the three Laurent Selmer hypotheses. The ideal $J$ is $S.J2$. A bound on a function into $F$ means that one power of
$t$ makes it integral and adically continuous.

% BEGIN P71S1
\paragraph{P71S1. The functionals on the two coefficient modules.}
The tangent space $T$ consists of $O'$-derivations $R\to F$
vanishing on $I$, and $L$ consists of additive $\mu:N\to F$
with $\mu(by)=r(b)\mu(y)$. There are identifications
\[
 T=(\mathfrak m/J)^*,\qquad L=(N/\mathfrak mN)^* .
\]
For the first, lift each residue $r(b)$ to $o\in O'$, and apply
a functional to the class of $b-o$. Two lifts differ by the
maximal ideal of $O'$, and the product rule follows by expanding
modulo $\mathfrak m^2$ as in CRW7. Vanishing on $I$ accounts
for the extra summand here. The second identification follows
from the scalar rule. These spaces are finite-dimensional by
Noetherianity of $R$ and finiteness of $N$. In particular every
functional on $T$ is evaluation at some $X\in\mathfrak m$.

For $\Phi:N\to F$ additive, $O'$-linear via the residue, and
zero on $JN$, put
\[
 D_y(b)=\Phi(by)-r(b)\Phi(y).
\]
Then $D_y\in T$, and $D_{b y}=r(b)D_y$ as a functional.
Indeed write $b=o+x$ using surjectivity from $O'$, and expand
products in two such expressions. The terms in $\mathfrak m^2N$
vanish; the scalar terms are taken outside $\Phi$. The same
calculation gives zero on $O'$ and $I$. Therefore for a basis
$\alpha_d$ of $T$ there are $\mu_d\in L$ such that
$D_y=\sum_d\mu_d(y)\alpha_d$. If $T=0$ this sum is empty.
Conversely those $\Phi$ separate $N/JN$, since $N/JN$ is a
vector space over $O'/\mathfrak m_{O'}=F$. This last test is
on $JN$, which may be smaller than $\mathfrak m N$.
% END P71S1

% BEGIN P71S2
\paragraph{P71S2. Bounded deformation tests preserving the ordered flags.}
Use the algebra with basis $1,e_d,f_d,\eta$ and multiplication
$e_d f_h=\delta_{dh}\eta$ as in Q70C2. Given $\alpha_d\in T$,
$\mu_d\in L$ and $\Phi$ with the product identity of P71S1, put
\[
 \theta(b+y)=r(b)+\sum_d f_d\alpha_d(b)+\sum_d e_d\mu_d(y)
                    +\eta\Phi(y) \qquad(b\in R,~y\in N).
\]
This is a ring map. Products in $R$ use the derivation law,
products between $R$ and $N$ use the identity for $\Phi$, and
both sides vanish on products in $N^2$ because $r$ and
$\alpha_d$ kill $I$. On $O'$ the map is the given scalar map.
Its application to $v$ has even entries in $1,f_d$ and odd
entries in $e_d,\eta$. It has the prescribed polarization,
determinant, inertia and ordered ordinary product relations:
the input representation has the form and determinant identities
and kills the defining group-algebra ideal, and the images of
the universal ordered local characters still supply the scalars
in the product. For the torsion subgroup their prescribed Teichmüller
values map to their constant residues because $O$ maps through
$O'$ in $R$, giving value one by $x.\psi_\mu$.

All coefficients of this representation are bounded. To verify it,
write $f:B\twoheadrightarrow A$ for the point and
$\mathfrak p=\ker f$. For each $b\in B$ the source provides
$b_0\in B$ with $f(b_0)=f(b)$ whose image in $R$ comes from
$O'$. Thus $\theta(b_0)=f(b)$. The Noetherian ideal
$\mathfrak p$ has finitely many generators $p_j$. Expand
$b-b_0$ in those generators and repeat on the coefficients.
Their images under $\theta$ are in the nilpotent augmentation
ideal of cube zero. Consequently $\theta(B)$ is contained in
the $A$-span of the finitely many monomials of degree at most two
in $\theta(p_j)$. Since $f$ is a local surjection,
$\theta(\mathfrak m_B^s)$ tends uniformly to zero: in a product of
$s$ factors at least $s-2$ contribute a residue in $tA$. The
other factors belong to the fixed bounded lattice just obtained.
Now use the module $S.U.P$, a finitely generated
$B$-module of matrices containing all $v(\gamma)$. Its image
under $\theta$ is bounded by the same argument on its generators.
The congruence condition on $v(\gamma u)-v(\gamma)$, modulo
$\mathfrak m_B^s S.U.P$ on an open normal subgroup, proves
adic continuity after clearing a uniform denominator. The same
bounds apply after multiplying by $\rho^{-1}$.

We record precisely how to use the ordered relation in these
tests. Compare two representations over this coefficient algebra
agreeing modulo a square-zero ideal $H$, with the first already
upper triangular in its ordered local basis and with ordered
characters distinct after reduction to $F$. Assume both satisfy
the vanishing ordered product and the two systems of characters
agree modulo $H$. Choose a group-algebra element distinguishing
all reduced characters, by the finite-hyperplane argument of CRW9.
Lagrange interpolation gives complete ordered idempotents for both
operators. The ordered relations with one varying argument give
zero from line $j$ to line $i$ when $i>j$ and the specified scalar
on each diagonal line, just as in CRW9. The intertwining sum of
the new projectors times the old ones is $1+U$ with entries in
$H$. It lies in the matrix subalgebra of the tests, by the polynomial
construction of its projectors and scalar denominators (their
inverses have the same allowed diagonal-block pattern).
Thus conjugation restores the ordered flag, with a correction in
the intersection of that subalgebra with $M_n(H)$. Its diagonal
comparison gives zero variation of the corrected cocycle on the
torsion subgroup. This argument also applies when the first
representation has nonconstant nilpotent entries: interpolation
only divides differences with nonzero residue, and its triangular
idempotents preserve the original ordered flag.
% END P71S2

% BEGIN P71S3
\paragraph{P71S3. The tangent and odd coefficient maps are injective on classes.}
Use the direct summands of $v(\gamma)$ entrywise to define
\[
 z_\alpha(\gamma)=\alpha(v_R(\gamma))\rho(\gamma)^{-1},
 \qquad u_\mu(\gamma)=\mu(v_N(\gamma))\rho(\gamma)^{-1}.
\]
The tests of P71S2 with a single $f$-direction or a single
$e$-direction and all other terms zero show these are bounded
even and odd Selmer cocycles, respectively. Multiplicativity
gives the cocycle equation. The form gives negative transpose
at conjugation, and the inertia equations give zero on prescribed
inertia. For the first test use the even square-zero corrections
of P71S2, and for the second use its odd corrections. These put
the local cocycles in the ordered flags with the required diagonal
on the torsion subgroup. Both maps are $F$-linear.

Suppose $z_\alpha$ is a skew even coboundary. Conjugate its
square-zero test to the constant $\rho$ by $1+fY$, with $Y$
its primitive. The determinant law implies $\alpha$ vanishes
on the determinant generators. To recover the ordered characters,
in the mixed ordinary product place a given element in slot $i$
and put in each other slot $j$ an element discriminating $i,j$.
The $i$th diagonal entry, now computed on the constant triangular
representation, forces the variation in slot $i$ to vanish since
all other factors are units. The auxiliary $U_{ji}$ are detected
by applying the first auxiliary polynomial relation to an eigenvector
of the corresponding reduced Frobenius with that labeled eigenvalue.
The distinct other factors are units. The second relation on the
same eigenvector forces the variation of $y_{ji}$ to vanish because
auxiliary inertia specializes to one. Both polynomial relations
hold on $v$ by the defining kernel condition, as in CRW11.
The parameter and $O$ come from $O'$ and are killed. These are
all generators in $S.P.\mathrm{gen}$. The continuity established in
P71S2 gives zero on $B$, hence on $B_{\mathfrak p}$ by the
inverse rule, and on $R$ since this localization surjects onto
the completion modulo $\mathfrak m^2$. Thus $\alpha=0$.

If $u_\mu$ is a coboundary, multiplying its equation on the
right by $\rho(\gamma)$ gives
$\mu(v_N(\gamma))=[\rho(\gamma),Y]$. The identity is
$R$-linear using residual scalars, so it holds throughout the
given span of $v$, with $\rho$ replaced there by $r(v_R)$.
Evaluate at a single off-block matrix unit with arbitrary
coefficient $y\in N$. Its residual matrix is zero, so
$\mu(y)=0$. If the second block is empty, the given type-I
field already gives $N=0$. Thus the map on $L$ is injective
modulo odd coboundaries in all cases.

The pulled-back $\kappa_j$ are independent in odd Laurent cohomology.
Indeed a Laurent relation downstairs on the source of the surjection
$S.h$ descends as an identity using the same primitive on the base
group, then multiplication by one power of $t$ makes it integral;
the original $\mathrm{Kap.indep}$ applies. By $K_3$ each class of
$u_\mu$ has unique coordinates in these generators. These
coordinates give an injection $L\hookrightarrow F^g$ since
$[u_\mu]=0$ forces $\mu=0$. By dualizing and using P71S1,
choose $m_j\in N$ lifting the images of the coordinate vectors
in $N/\mathfrak mN$. They span $N$ by Nakayama, and
\[
       [u_\mu]=\sum_j\mu(m_j)[\kappa_j]
                  \qquad(\mu\in L).
\]
% END P71S3

% BEGIN P71S4
\paragraph{P71S4. A lift attached to each functional on $N/JN$.}
Take $\Phi$ as in P71S1, write $D_y=\sum_d\mu_d(y)\alpha_d$,
and put $C_{dj}=\mu_d(m_j)$. Apply the test map of P71S2 to $v$.
Its normalized $e_d$ and $f_d$ coefficients, obtained by multiplying
on the right by $\rho(\gamma)^{-1}$, are respectively
$u_{\mu_d}$ and $z_{\alpha_d}$. By P71S3 the first differs
from $\sum_j C_{dj}\kappa_j$ by a coboundary with skew odd
primitive $Y_d$. Conjugating by $1+\sum e_dY_d$ removes this
difference, since the change to the $e_d$ coefficient is
$Y_d-\rho Y_d\rho^{-1}$. It preserves the $f_d$ coefficients,
the ordered relations and the form, the conjugating matrix being
orthogonal.

Change to an algebra with $g$ pairs, using the following relation
between old and new coefficients. The constant and $\eta$
coordinates agree, the old $e$-column is $C$ times the new
$e$-column, and the new $f$-column is $C^T$ times the old
$f$-column. Related pairs are closed under addition and
multiplication: the contractions $e\cdot f$ agree by the
transpose rule. There is no invertibility condition on $C$.
Prescribe the new $e_j$ coordinate of the representation to be
$\kappa_j(\gamma)\rho(\gamma)$, and of each local character
to be zero. This produces related pairs by the preceding
conjugation. In the quotient by the square-zero ideal $(f_j,\eta)$
the new representation is the original
$R_1=(1+\sum e_j\kappa_j)\rho$ and its characters are the
original $\psi_i$. Hence its multiplication, form, inertia, and
ordered relations hold in that quotient by T70B1. For each such
identity the constant, $\eta$ and $f$ coordinates additionally
agree by the relation to the old identity. Together these two
checks verify all coordinates. Boundedness follows entrywise
from P71S2, the fixed finite operations, and the original
integral $\kappa_j$.

Denote the new representation by $U$. Its difference
$Z=U R_1^{-1}-1$ has only $f_j$ on the even blocks and $\eta$
on the odd blocks, the exact space $W_2$. The coefficient ideal
of this space is square-zero. Multiplication thus gives the
$R_1$-cocycle relation, polarization gives the required negative
transpose, and inertia gives zero. Apply the square-zero flag
argument of P71S2 comparing $U$ with $R_1$ in its deformed
ordered basis. The correcting matrix belongs to $W_2$: it lies
in $(f_j,\eta)$ and in the matrix subalgebra with even entries
in $1,f_j$ and odd entries in $e_j,\eta$. Thus $Z$ satisfies
the original $\mathrm{Sel2}$ conditions including the torsion
diagonal, and it is bounded. Its $f_j$ coefficient is
$z_{D_{m_j}}$, by linearity and the transpose rule. We have
proved that $(D_{m_j})_j$ belongs to
\[
 \mathcal U=\{(\alpha_j)_j\in T^g:
            (z_{\alpha_j})_j\text{ admits a bounded Sel2 lift}\}.
\]
This is a vector subspace since one can add lifts with their
local corrections and scale by any fixed element of $F$.
% END P71S4

% BEGIN P71S5
\paragraph{P71S5. Coordinatewise lifts give a first-order cocycle.}
Let $z$ be a bounded even Selmer cocycle such that for each $j$
the tuple with $z$ in slot $j$ and zero elsewhere has a bounded
$W_2$-lift $Z_j$. Write $H_j=[\eta]Z_j$ and put
$y=z+\sum_j e_jH_j$. This is $W_1$-valued, and multiplication
by $f_j$ gives exactly $Z_j$. An element of $W_1$ is determined
by its constant coefficient and all its products by $f_j$.
Thus the cocycle identity, the sign and the inertia condition
follow by checking on $z$ and the $Z_j$ (in the constant term
$R_1$ acts as $\rho$). Boundedness and the required continuity
hold coordinatewise. For $g=0$ only the constant test is needed.

Here are compatible local witnesses. Fix an even correction $Y$
for $z$ and $W_2$ corrections $Y_j$ for $Z_j$. For each $i,j$
let $Q_{ji}=[f_i]Y_j-\delta_{ij}Y$, an even matrix. Its local
coboundary has zero entries below the diagonal in the $w$ basis
by subtracting the two triangularity conditions in the $f_i$
coordinate. Thus the class of $Q_{ji}$ modulo upper triangular
matrices is a local invariant. The strictly lower quotient has
a filtration with characters $\psi_p/\psi_q$, $p>q$, by the
entry calculation in U70O5, now on even entries. By $x.N1$ none
is trivial over $F$. Hence $w^{-1}Q_{ji}w$ is upper triangular.

Write $E_v=\sum_i e_i S.k.Y_{i,v}$ and
$W=(1-E_v)w$. Subtract from $Y_j$ the matrix
\[
    (1-E_v)\left(\sum_i f_i Q_{ji}\right)(1+E_v).
\]
This subtraction preserves the permitted corrections: the
subtracted term is in $W_2$, is upper triangular in the $W$
basis, and its coboundary there is upper triangular with zero
diagonal, by triangularity of $R_1$. The adjusted $Y_j$ have
$f_i$ coefficients $\delta_{ij}Y$. Their $\eta$ coefficients
$H'_j$ now give a simultaneous $W_1$ correction
$Y'=Y+\sum e_j H'_j$, since $f_jY'$ equals the adjusted $Y_j$.
Testing the desired corrected lower and torsion-diagonal entries
on their constant coefficients and on their products by $f_j$
proves both vanishings. We have constructed $y\in\mathrm{Sel1}$
with constant part $z$.

Define $W_0=\{\alpha\in T:\text{each one-slot tuple of }\alpha
\text{ lies in }\mathcal U\}$. For a basis
$\alpha_1,\ldots,\alpha_h$ of $W_0$, apply the construction to
each $z_{\alpha_i}$. Add to these $h$ lifts the $g^2$ cocycles
$e_j\kappa_i$ in $W_1$. They are bounded with the required
local corrections $e_j S.k.Y_{i,v}$ because multiplication by
$e_j$ kills the $e$-terms in $R_1$ and in $W$. These $h+g^2$
classes are independent. Indeed a relation modulo a skew
$W_1$-coboundary with primitive $Y$ gives, in its constant
coordinate, a relation of the $z_{\alpha_i}$ modulo an even
boundary. By P71S3 their coefficients vanish. The remaining
constant primitive is skew and commutes with $\rho$. The span
of the cyclotomic kernel is the product of matrix algebras on
the nonempty blocks, so a commuting even matrix is scalar on
each. Skewness forces zero in odd characteristic. In the
$e_j$ coordinate one is left with a relation among the pulled-back
$\kappa_i$ modulo an odd skew coboundary. P71S3 forces all
remaining coefficients zero. Now $K_1(D)$ gives
\[
                         \dim W_0+g^2\le D .
\]
% END P71S5

% BEGIN P71S6
\paragraph{P71S6. The independent coefficient rows and their relations in $N$.}
By P71S3 every tuple in $T^g$ gives a tuple of bounded even
Selmer cocycles. The hypothesis $K_2$ says that every $g+1$
elements of $T^g$ are dependent modulo $\mathcal U$. Therefore
$\dim(T^g/\mathcal U)\le g$. Choose $g$ linear forms spanning
the annihilator of $\mathcal U$, padding with zeros, and decompose
them as $\Lambda_i(\alpha)=\sum_j\lambda_{ij}(\alpha_j)$
with $\lambda_{ij}\in T^*$. The common kernel of the
$\lambda_{ij}$ inside $T$ is exactly $W_0$, by testing
each one-slot vector against all $\Lambda_i$. Their span in
$T^*$ therefore has codimension $\dim W_0$. By P71S5 we can
form a spanning list of length $D$ in $T^*$ by putting the
$g^2$ individual forms $\lambda_{ij}$ in distinct slots
$p(i,j)$, appending a basis of a complement, then appending
zero forms if required. The functionals themselves may coincide
or vanish; the distinct slots give $g^2$ independent coefficient
rows. For $g=0$ the same argument just chooses a list spanning
$T^*$.

Lift this list to elements $X_a\in\mathfrak m$ by P71S1.
They span $\mathfrak m/J$ over $F$ and hence
$\mathfrak m\subset (X_a)+J$ over $R$. Use the $m_j$ of
P71S3 and take $c_{ij,a}\in O'$ to be one if $a=p(i,j)$ and
zero otherwise. Their reductions are independent coordinate
rows in $(O'/\mathfrak m_{O'})^D$. For any $\Phi$ on $N/JN$,
P71S4 gives $(D_{m_j})\in\mathcal U$. Thus for each $i$
\[
  0=\sum_j\lambda_{ij}(D_{m_j})
     =\sum_j D_{m_j}(X_{p(i,j)})
     =\Phi\!\left(\sum_j X_{p(i,j)}m_j\right),
\]
since $r(X_a)=0$. The separation in P71S1 proves the last
sum lies in $JN$. It is exactly
$\sum_j(\sum_a c_{ij,a}X_a)m_j$ in the native definition,
with the $O'$ coefficients viewed in $R$. Together with the
ideal containment and the module generators already obtained,
these data give every field of $S.\mathrm{Pres}(D)$.
This proves $\mathrm{Pa.pres}$ from the three original bounded
Laurent Selmer hypotheses at arbitrary $g$.
% END P71S6
